% Sources:
% - manuscript/sections/02-batchreceive-identity-problem.tex
% - manuscript/sections/03-security-objective-threat-model.tex
% - manuscript/sections/04-formal-modeling.tex
% - manuscript/sections/05-formal-analysis.tex
% - docs/paper/discussion-draft.md
% - docs/paper/limitations-draft.md
% - docs/paper/contribution-evidence-map.md
% - docs/rq-v2/rq-v2-complete-argument-map.md
% - docs/rq-v2/g2-admission-semantics.md
% - docs/rq-v2/rq-v2-threat-model.md
% Supporting invariant/enforcement and interface analyses were consulted only
% for interpretation; they supply no additional machine-checked result.

\subsection{Why the Result Is More Than Guard Removal}

Removing a condition \(A_1\neq A_2\) unsurprisingly permits the variables to
be equal. A comparison that stopped at this syntactic observation would say
little about the role of the condition. The question addressed by the
controlled analysis is instead which identity relation must be preserved to
maintain the intended interpretation of one admitted batch.

Section~5 connects this question to three complementary observations. In the
relaxed model, repeated input participates in a complete reachable lifecycle:
one matching modeled \(\mathsf{Send}\) origin precedes one admitted batch and
two distinct \receiveraccept occurrences for the same tuple and batch
context. This connects structural relaxation to receiver-event multiplicity,
rather than merely to a possible assignment of party variables. The
observable endpoint remains a modeled receiver event, not an application
action.

The message-level variant supplies the decisive control. It excludes the
exact duplicate-acceptance pattern and satisfies message distinction and
scoped occurrence injectivity, yet permits two different messages from the
same party to be accepted together. Those acceptances have different modeled
sender origins. Thus, success at preventing repeated acceptance of one exact
origin does not resolve the party-composition question. Finally, the
party-level variant preserves the target relation while a valid distinct-party
batch remains reachable; the resulting safety property is therefore
non-vacuous.

The comparison in Figure~\ref{fig:identity-control-comparison} consequently
distinguishes the target semantic relation from plausible substitutes over
other coordinates. It is a protocol-specific analysis of the stated
\batchreceive condition, not a new general verification methodology. Its
value lies in identifying what the modeled controls preserve and what their
success leaves undecided.

\subsection{Party Identity, Message Identity, and Occurrence Injectivity}

The three dimensions answer different questions. Party identity determines
which modeled principal an entry projects to. Message identity compares its
symbolic message coordinate with that of another entry. Scoped occurrence
injectivity concerns whether one matching sender occurrence can support more
than one matching acceptance occurrence in the same batch and receiver
context. The last relation is about event multiplicity, not the number of
distinct principals represented by different messages.

One party \(A\) may originate tuples with
\(\mathit{sid}_1\neq\mathit{sid}_2\) and \(m_1\neq m_2\). Accepting each of
these tuples once does not repeat either exact sender origin, but placing
both in one batch repeats the party projection. Equation~\eqref{eq:message-not-party}
expresses the specific non-implication exposed by the message-level control:
message distinction together with scoped occurrence injectivity does not
entail party distinction.

These properties should therefore not be read as successive guarantees in a
hierarchy culminating in same-batch party distinction. Nor does this comparison establish
every possible logical independence among them; relationships can depend on
the modeled freshness, matching, and scope assumptions. Occurrence
injectivity serves here as a diagnostic and comparison property. It separates
reuse of one exact origin from legitimate multiple origins belonging to one
party. The primary objective remains the party-level composition invariant,
not injectivity as a newly introduced notion.

\subsection{Invariant and Enforcement}

\distinctparty states what relationship holds among entries of one admitted
batch, as defined in Equation~\eqref{eq:distinct-party-objective}. An
enforcement mechanism determines how that relationship is maintained. The
party-level rule directly compares \(A_1\) and \(A_2\), providing one
executable realization of the invariant in the symbolic abstraction. This
realization makes the comparison concrete; it does not turn that particular
check into the only permissible implementation.

In principle, a trusted caller could construct only party-distinct batches,
a batch builder could maintain the relation during composition, or a
receiver-side admission layer could check it before processing. Another
component with the relevant party attribution could also be responsible.
These are architectural possibilities, not evaluated alternatives or findings
about where a deployed \kwaay implementation performs enforcement.

The design obligation is conditional on the intended semantics: if different
batch positions are interpreted as contributions from different parties, the
component responsible for construction or admission must preserve that
party-level relationship. Necessity here concerns preserving this
interpretation at the modeled boundary, not every cryptographic property of
\kwaay or one mandatory checking algorithm.

The representation used for that obligation also needs to be explicit. The
symbol \(A\) is a protocol-principal coordinate, not automatically an account
identifier, public-key encoding, username, or database record. A concrete
integration needs a justified mapping to that coordinate before comparing
representations. Neither that mapping nor the trust placed in the component
that supplies it is validated by the present model.

\subsection{Relation to Authentication, Replay, and Deduplication}

An entry-level origin relation and a relation among batch participants are
different obligations. In the current abstraction, sender-origin
correspondence says that an accepted tuple has a prior matching
\(\mathsf{Send}(A,\mathit{sid},m)\). It does not establish full \kwaay
authentication. Even this correspondence can hold for both entries while
their joint composition repeats a party, as the message-level example
demonstrates. Entry validity, in this modeled origin sense, does not imply
composition validity. More generally, an authentication contract that only
relates each item to its origin does not additionally state pairwise party
inequality across batch positions.

The relaxed witness has a replay interpretation because the same exposed
tuple is supplied to two positions. The identity question is nevertheless
not exhausted by that pattern. The message-level witness uses different
messages and sender occurrences, so its party repetition does not depend on
replaying one exact message. Removing the first pattern therefore leaves a
separate composition case to consider.

Exact-message deduplication succeeds at the purpose encoded in the control:
the equal-message rejection branch is reachable, accepted messages are
distinct within the modeled context, and the scoped occurrence-injectivity
property holds. Its insufficiency for same-batch party distinction is not evidence that
deduplication is useless. It reflects a difference between the coordinate
being compared and the relation being sought. Individual-entry properties,
exact-repeat controls, and batch-composition properties consequently warrant
separate analysis rather than treatment as interchangeable solutions. Here
that comparison uses fresh symbolic messages, which associate distinct sender
occurrences with distinct message coordinates. It is not an evaluation of a
wire-format deduplication algorithm or a cross-batch replay cache.

\subsection{Design Implications}

\paragraph{Treat composition as a semantic boundary.}
When several individually legitimate entries are combined into one protocol
action, their joint arrangement can be subject to an additional condition.
For the interface studied here, checking origins does not discharge the
distinct-party obligation. A specification or integration contract should
state the admissible participant relationship as well as the conditions on
individual entries.

\paragraph{Constrain the coordinate used by the interpretation.}
If the required interpretation is distinct parties, inequality of messages,
sender sessions, or slot indices cannot stand in for party inequality merely
because each distinguishes some aspect of an entry. A proposed substitute
needs an argument connecting its representation and comparison to the target
relation. The message-level control supplies a concrete case where that
connection does not hold.

\paragraph{Assign responsibility for preserving the relation.}
An integration should make explicit which component maintains the stated
condition and what party attribution that component relies on. The condition
reported in Section~2 is not evidence that a deployed implementation omits
enforcement, and the present analysis does not prescribe adding a particular
check to \kwaay. It identifies an obligation whose realization may belong to
the caller, batch builder, or admission layer.

Together, these implications suggest a discipline for this interface:
identify the semantic coordinate on which composition depends, state its
invariant, and assign responsibility for preserving it. Any propagation from
duplicate acceptance to application state would additionally require an
explicit consumer model; it is outside the direct evidence presented here.

\subsection{Limitations}

\paragraph{Fixed two-slot evidence.}
The semantic definition allows arbitrary finite batch size, but the
machine-checked evidence fixes \(n=2\). Two slots represent the smallest
violating pair; this does not establish how larger batches behave or prove
the comparison for arbitrary \(n\). The claims are batch-local and do not
supply cross-batch, rollback, or restart guarantees.

\paragraph{Admission-level abstraction.}
The model isolates composition, admission, and origin-matched receiver
events rather than the complete symbolic \kwaay construction. KEM and
signature computations, prekey cryptography, key derivation, ratcheting, and
compromise behavior are omitted. Fresh sender-session and message coordinates
and persistent origin facts are explicit prototype choices, not a
reconstruction of all concrete message behavior. The analysis therefore
provides neither a full-protocol security proof nor evidence of a
cryptographic secrecy or authentication failure. Table~\ref{tab:source-abstraction}
separates source-derived roles from these assumptions; no refinement theorem
connects the abstract receiver events to concrete output keys. Adversarial
control over composition is an analysis assumption, not an established
capability against a deployed batch builder.

\paragraph{Receiver-event endpoint.}
Direct evidence ends at \receiveraccept. The prototypes contain no output-key
objects, \textsc{Key}/\textsc{Test} interfaces, or application consumer.
Duplicate acceptance alone consequently establishes neither failure of a
key interface nor duplicate installation, state corruption, or application
compromise. Conceptual motivations involving downstream interpretation do
not extend this evidence boundary.

\paragraph{Rejection evidence.}
The rejection properties establish existence of a trace, not eventual
rejection of every invalid collected batch. Moreover,
\path{same_party_rejection_exists} does not bind its two \(\mathsf{Send}\)
tuples to the rejected batch. As distinguished in Section~5, rejection of a
collected same-party/different-message pair follows from the rule semantics,
not from a tuple-bound witness supplied by that lemma.
The prototype records an abstract batch-rejected state. It does not establish
whether a concrete integration would discard one entry, reject the whole
batch, return an error, or expose another failure behavior.

\paragraph{Representation and enforcement.}
The model treats party coordinates as explicit symbolic values. It does not
validate identity resolution, key-to-party or database mappings, enforcement
by a caller, or distributed maintenance of the relation. The proposed
placement options are therefore design considerations rather than verified
implementation choices or evidence of a deployed defect.

Finally, the artifact records model hashes, commands, environment, and
outcomes, but contains no raw prover transcripts; manuscript preparation did
not include an independent prover rerun. Larger-batch analyses and explicit
protocol or consumer refinements would be needed to assess extensions beyond
the present boundary.
